\begin{MyRecipe}{Joscha's Gefüllte Paprika}{\Calc{4}{\x} Personen}{\SI{30}{\minuteprime} + \SI{90}{\minuteprime} }
	

	\ingredient[\Calc{1}{\x}]{\si{\kilogram}} {Tomaten}
	\ingredient[\Calc{100}{\x}]{\si{\gram}} {Tomatenmark}
	\ingredient[\Calc{2}{\x}]{\si{\liter}} {Wasser}
	\ingredient[]{} {Salz}
	\ingredient[]{} {Zucker}
 	\ingredient[]{} {Pfeffer}

	\Step{Tomaten passieren}
	Tomaten kochen und durchpassieren. Brühe aus passierten Tomaten, Tomatenmark, Wasser  und Salz verquirlen. 

	\ingredient[\Calc{0.5}{\x}]{\si{\kilogram}} {Hackfleisch, gem.}
	\ingredient[\Calc{2}{\x}]{} {Peperoni}
	\ingredient[\Calc{2}{\x}]{} {Zwiebeln}
	\ingredient[\Calc{2}{\x}]{} {Knoblauchzehen}
	\ingredient[\Calc{1}{\x}]{} {Ei}
	\ingredient[\Calc{160}{\x}]{\si{\gram}} {Reis}

	\Step{Füllung}
	Zwiebeln, Knoblauch und Peperoni fein hacken.

	Teig aus Fleisch, Gehacktem, Reis, Ei, Salz und Pfeffer kneten.\par\bigskip

	\ingredient[\Calc{12}{\x}]{} {Spitzpaprika}

	\Step{Paprikas füllen}
	Spitzpaprika vorsichtig vom Strunk her aushöhlen und waschen und Füllung in die Paprikaschoten verteilen.

	\ingredient[\Calc{3}{\x}]{} {Schweinsknöchle}
	\ingredient[\Calc{2}{\x}]{} {Schweinsfüße}
	\ingredient[\Calc{1}{\x}]{\si{\Essloeffel}} {Schweineschmalz}

	\Step{Zubereitung}
	Schweinsknöchle und -füße mit einem Messer gründlich abschaben und waschen. In Schweineschmalz anbraten und mit der Brühe ablöschen.
	
	Gefüllte Paprika einlegen und mit Wasser aufgießen, bis die Paprika fast bedeckt sind.
	
	\SI{90}{\minuteprime} köcheln lassen.	\par\bigskip


	\Step{\emph{Hinweise}}
	\begin{itemize}
		\item Anstelle der selbst passierten Tomaten können auch einfach \Calc{1}{\x}\,\si{\kilogram} passierte Tomaten aus dem Supermarkt verwendet werden. 
	
		\item Anstelle der Peperoni kann auch mit Harissa nach bedarf gewürzt werden.
	\end{itemize}
	\par\bigskip



	
	
	
\end{MyRecipe}